\documentclass[12pt]{article}
\usepackage[UTF8]{inputenc}
\usepackage{thaisupp}
\usepackage{enumerate}
\usepackage{etoolbox}
\AtBeginEnvironment{enumerate}{\itshape}
\begin{document}
\title{การเคลื่อนที่แบบโปรเจกไทล์}
\maketitle
\section*{คำถาม in class}
\begin{enumerate}
\def\labelenumi{\arabic*.}
\normalfont
\item \textbf{วัตถุถูกขว้างออกไปในแนวราบด้วยความเร็ว 20 m/s จากหน้าผาสูง 80 m ถามว่าความเร็วในแนวดิ่งเมื่อถึงพื้นเป็นเท่าใด (g = 10 m/s²)}
\begin{enumerate}[label=\alph*)]
\item 20 m/s
\item 30 m/s
\item 40 m/s
\item 50 m/s
\end{enumerate}
\item \textbf{ขว้างวัตถุในแนวราบด้วยความเร็ว 15 m/s จากหน้าผาสูง 45 m วัตถุตกถึงพื้นที่ระยะทางแนวราบเท่าใด (g = 10 m/s²)}
\begin{enumerate}[label=\alph*)]
\item 30 m
\item 45 m
\item 60 m
\item 75 m
\end{enumerate}
\item \textbf{ยิงวัตถุขึ้นไปในแนวมุม 60° กับแนวราบด้วยความเร็ว 40 m/s วัตถุใช้เวลาขึ้น-ลงทั้งหมดเท่าใด (g = 10 m/s²)}
\begin{enumerate}[label=\alph*)]
\item 2 s
\item 4 s
\item 6 s
\item 8 s
\end{enumerate}
\item \textbf{การเคลื่อนที่แบบโปรเจกไทล์ วิถีการเคลื่อนที่ (trajectory) มีรูปร่างเป็นอะไร}
\begin{enumerate}[label=\alph*)]
\item เส้นตรง
\item วงกลม
\item พาราโบลา
\item วงรี
\end{enumerate}
\item \textbf{ในการเคลื่อนที่แบบโปรเจกไทล์ ข้อใดต่อไปนี้ถูกต้อง}
\begin{enumerate}[label=\alph*)]
\item ความเร็วในแนวราบเปลี่ยนแปลงตลอดเวลา
\item ความเร็วในแนวดิ่งเปลี่ยนแปลงตลอดเวลา
\item ความเร็วในแนวราบและแนวดิ่งเปลี่ยนแปลงเท่ากัน
\item ความเร็วในแนวดิ่งคงที่
\end{enumerate}
\item \textbf{ยิงวัตถุขึ้นไปในแนวมุม 30° กับแนวราบด้วยความเร็ว 60 m/s ความสูงสูงสุดที่วัตถุจะไปถึงเท่ากับ (g = 10 m/s²)}
\begin{enumerate}[label=\alph*)]
\item 45 m
\item 90 m
\item 135 m
\item 180 m
\end{enumerate}
\item \textbf{ยิงวัตถุในแนวมุม 45° กับแนวราบด้วยความเร็ว 40 m/s ระยะไกลสุดในแนวราบเท่ากับ (g = 10 m/s²)}
\begin{enumerate}[label=\alph*)]
\item 80 m
\item 120 m
\item 160 m
\item 200 m
\end{enumerate}
\item \textbf{วัตถุถูกขว้างในแนวมุม 37° กับแนวราบด้วยความเร็ว 50 m/s ขนาดของความเร็วลัพธ์เมื่อวัตถุอยู่ที่จุดสูงสุดเท่ากับ (g = 10 m/s², sin 37° = 0.6, cos 37° = 0.8)}
\begin{enumerate}[label=\alph*)]
\item 30 m/s
\item 40 m/s
\item 50 m/s
\item 60 m/s
\end{enumerate}
\item \textbf{ขว้างวัตถุจากหน้าผาสูง 180 m ขึ้นไปในแนวมุม 30° กับแนวราบ วัตถุใช้เวลากี่วินาทีกว่าจะตกถึงพื้น (g = 10 m/s²)}
\begin{enumerate}[label=\alph*)]
\item 6 s
\item 8 s
\item 10 s
\item 12 s
\end{enumerate}
\item \textbf{ถ้าต้องการยิงวัตถุให้ได้ระยะไกลสุดในแนวราบบนพื้นราบ ควรยิงที่มุมเท่าใด (ไม่คิดแรงต้านอากาศ)}
\begin{enumerate}[label=\alph*)]
\item 30°
\item 45°
\item 60°
\item 90°
\end{enumerate}
\item \textbf{เครื่องบินบินในแนวราบด้วยความเร็ว 200 m/s ที่ความสูง 500 m ปล่อยลูกน้ำหนักลงพื้น ลูกน้ำหนักจะตกห่างจากจุดปล่อยกี่เมตร (g = 10 m/s²)}
\begin{enumerate}[label=\alph*)]
\item 1,000 m
\item 1,414 m
\item 2,000 m
\item 2,828 m
\end{enumerate}
\item \textbf{ในการเคลื่อนที่แบบโปรเจกไทล์ ถ้าไม่คิดแรงต้านอากาศ ข้อใดต่อไปนี้ถูกต้องเกี่ยวกับการเคลื่อนที่ในแนวราบ}
\begin{enumerate}[label=\alph*)]
\item เป็นการเคลื่อนที่ด้วยความเร่งคงที่
\item เป็นการเคลื่อนที่ด้วยความเร็วคงที่
\item เป็นการเคลื่อนที่ด้วยความเร่งเปลี่ยนแปลง
\item วัตถุหยุดนิ่งในแนวราบ
\end{enumerate}
\item \textbf{ยิงลูกบอลขึ้นไปในแนวมุม 53° กับแนวราบด้วยความเร็ว 50 m/s (sin 53° = 0.8, cos 53° = 0.6) ความสูงสูงสุดเท่ากับ (g = 10 m/s²)}
\begin{enumerate}[label=\alph*)]
\item 40 m
\item 60 m
\item 80 m
\item 100 m
\end{enumerate}
\item \textbf{ยิงวัตถุที่มุม 30° กับแนวราบด้วยความเร็ว 30 m/s จากหน้าผาสูง 20 m (วัตถุตกต่ำกว่าจุดยิง 20 m) ระยะไกลสุดในแนวราบเท่ากับ (g = 10 m/s²)}
\begin{enumerate}[label=\alph*)]
\item 26 m
\item 52 m
\item 78 m
\item 104 m
\end{enumerate}
\item \textbf{ยิงวัตถุด้วยความเร็วคงที่ v₀ ที่มุม θ และ (90° - θ) สองมุมนี้ให้ระยะไกลเท่ากันเมื่อ}
\begin{enumerate}[label=\alph*)]
\item θ = 30°
\item θ = 45°
\item θ = 60°
\item θ = 75°
\end{enumerate}
\item \textbf{ขว้างวัตถุในแนวมุม 45° กับแนวราบด้วยความเร็ว 20√2 m/s วัตถุมีขนาดความเร็วเมื่อถึงพื้นเท่ากับ (g = 10 m/s²)}
\begin{enumerate}[label=\alph*)]
\item 10 m/s
\item 20 m/s
\item 30 m/s
\item 40 m/s
\end{enumerate}
\item \textbf{ในการเคลื่อนที่แบบโปรเจกไทล์ ณ ตำแหน่งสูงสุดของวิถี ความเร็วของวัตถุมีค่าเท่าใด}
\begin{enumerate}[label=\alph*)]
\item 0 m/s
\item เท่ากับความเร็วเริ่มต้น
\item เท่ากับความเร็วในแนวราบเท่านั้น
\item เท่ากับความเร็วในแนวดิ่งเท่านั้น
\end{enumerate}
\item \textbf{ในการเคลื่อนที่แบบโปรเจกไทล์ ความเร่งของวัตถุมีค่าเท่าใด (ไม่คิดแรงต้านอากาศ)}
\begin{enumerate}[label=\alph*)]
\item 0 m/s²
\item 9.8 m/s² ในแนวราบ
\item 9.8 m/s² ในแนวดิ่ง
\item 9.8 m/s² ในทุกทิศทาง
\end{enumerate}
\item \textbf{นักกีฬาขว้างจักรยานอนในแนวมุม 45° ได้ระยะไกล 80 m ถ้าต้องการให้ได้ระยะไกล 100 m โดยยิงที่มุมเดิม ต้องเพิ่มความเร็วเริ่มต้นกี่เท่า}
\begin{enumerate}[label=\alph*)]
\item 1.00 เท่า
\item 1.25 เท่า
\item 1.50 เท่า
\item 1.75 เท่า
\end{enumerate}
\item \textbf{ยิงวัตถุในแนวมุม 53° ด้วยความเร็ว 50 m/s จากหน้าผาสูง 80 m (sin 53° = 0.8, cos 53° = 0.6) วัตถุอยู่สูงจากหน้าผากี่เมตรเมื่อถึงจุดสูงสุด (g = 10 m/s²)}
\begin{enumerate}[label=\alph*)]
\item 45 m
\item 80 m
\item 100 m
\item 125 m
\end{enumerate}
\item \textbf{วัตถุถูกปล่อยจากที่สูง h พร้อมกับวัตถุถูกขว้างในแนวราบจากที่สูงเดียวกัน วัตถุใดจะถึงพื้นก่อน (ไม่คิดแรงต้านอากาศ)}
\begin{enumerate}[label=\alph*)]
\item วัตถุที่ถูกปล่อย
\item วัตถุที่ถูกขว้าง
\item ถึงพื้นพร้อมกัน
\item ขึ้นกับมวลของวัตถุ
\end{enumerate}
\item \textbf{ยิงวัตถุในแนวมุม 60° กับแนวราบด้วยความเร็ว 40 m/s ความเร็วของวัตถุเมื่อผ่านจุดที่อยู่สูงจากพื้นเท่ากับจุดเริ่มต้น (กลับมาที่ระดับเดิม) มีค่าเท่าใด (g = 10 m/s²)}
\begin{enumerate}[label=\alph*)]
\item 20 m/s
\item 30 m/s
\item 40 m/s
\item 50 m/s
\end{enumerate}
\item \textbf{ขว้างวัตถุจากหน้าผาสูง 80 m ในแนวมุม 37° กับแนวราบ (sin 37° = 0.6, cos 37° = 0.8) วัตถุตกห่างจากฐานหน้าผา 120 m ความเร็วเริ่มต้นเป็นเท่าใด (g = 10 m/s²)}
\begin{enumerate}[label=\alph*)]
\item 20 m/s
\item 25 m/s
\item 30 m/s
\item 35 m/s
\end{enumerate}
\item \textbf{เครื่องบินทิ้งระเบิดบินขนานพื้นราบที่ความสูง 2,000 m ด้วยความเร็ว 360 km/hr ระยะห่างแนวราบจากเป้าที่ต้องทิ้งระเบิดเป็นเท่าใด (g = 10 m/s²)}
\begin{enumerate}[label=\alph*)]
\item 400 m
\item 600 m
\item 800 m
\item 1,000 m
\end{enumerate}
\item \textbf{ยิงวัตถุในแนวมุม 30° กับแนวราบจากหน้าผาสูง 100 m วัตถุตกถึงพื้นที่ระยะ 173 m จากฐานหน้าผา ความเร็วเริ่มต้นเป็นเท่าใด (g = 10 m/s², √3 ≈ 1.732)}
\begin{enumerate}[label=\alph*)]
\item 10 m/s
\item 20 m/s
\item 30 m/s
\item 40 m/s
\end{enumerate}
\end{enumerate}
\end{document}